\chapter{Conditioning the Problem}
\label{ch:conditioning}

Every iteration count in this book is governed by a condition number. CG needs
$O(\sqrt\kappa)$ iterations (Proposition~\ref{prop:cg}); projected gradient needs
$O(\kappa)$ (Proposition~\ref{prop:proj}); a direct solve loses digits in proportion
to $\kappa$. So anything that compresses the spectrum shortens the solve, and this
chapter is about the mechanisms that do.

There are three, and they are usually taught in three different courses. \emph{
Regularisation} moves the operator toward a target the application deems preferable.
\emph{Preconditioning} accelerates the solve of the operator as given, leaving the
answer untouched. \emph{Spectral cleaning} replaces the part of the spectrum that is
estimation noise. The chapter's thesis is that on the problems of this book the three
coincide in a useful way, and that a practitioner applying the first for statistical
reasons is getting the second for free---a connection that appears not to have been
made explicit in the application literature.

The running example is a sample covariance matrix, because it is the most
spectacularly ill-conditioned operator a reader is likely to meet: on five years of
S\&P~500 daily returns, $\kappa(\hat\Sig) \approx 10^5$.

\section{Regularisation}

Replace $H$ by the convex combination
\begin{equation}\label{eq:regsplit}
  H_\alpha = (1-\alpha)\,H + \alpha\,R\T R, \qquad \alpha \in (0,1),\; R\T R \succ 0,
\end{equation}
a ridge/Tikhonov regularisation toward the SPD target $R\T R$.
Proposition~\ref{prop:weyl} already gave the effect on the condition number, and for
the scaled-identity target it is $\kappa(H_\alpha) = O((1-\alpha)/\alpha)$: doubling
$\alpha$ roughly halves the bound, \emph{independently of how small the smallest
eigenvalue of $H$ was}.
\index{regularisation}

Two features of the split matter beyond the bound.

\paragraph{It is free in factored form.} When $H = M\T M$ the regularised operator is
again a Gram matrix, of the stacked operator
\[
  \tilde{M} = \begin{pmatrix} \sqrt{1-\alpha}\,M \\ \sqrt{\alpha}\,R \end{pmatrix},
  \qquad \tilde{M}\T\tilde{M} = H_\alpha,
\]
so it is simply the Gram backend of Chapter~\ref{ch:operator} applied to $\tilde M$,
and every solver runs on it with no change. The augmented-matrix device is standard in
ridge regression and Tikhonov regularisation \cite{golub2013,lawson1974}; what is worth
noticing is that it keeps the method \emph{matrix-free}.

\paragraph{It restores the $P$-matrix property.} Because $H_\alpha \succ 0$ strictly
whenever $\alpha > 0$, every principal submatrix is invertible even when $H$ is only
positive semidefinite---a rank-deficient least-squares operator with fewer rows than
columns, say. So the split simultaneously conditions the inner solve and
\emph{licenses the finite-termination proof} of Theorem~\ref{thm:termination}, which
needs the $P$-matrix property of Chapter~\ref{ch:geometry}. A strictly positive
$\alpha$ is therefore the natural setting for the active-set loop, and the two
motivations for it are independent.

\section{Preconditioning}

Regularisation lowers $\kappa$ by \emph{changing the operator}. When no modelling
change is admissible and $H$ must be solved as given, the same reduction in iteration
count is available through preconditioning, which leaves the minimiser untouched.

An SPD preconditioner $P_\Free \approx H_{\Free\Free}$ turns the inner solve into
preconditioned CG, equivalently CG on
$P_\Free^{-1/2} H_{\Free\Free} P_\Free^{-1/2}$, whose spectrum coincides with that of
$P_\Free^{-1}H_{\Free\Free}$. PCG realises this without ever forming
$P_\Free^{-1/2}$: each iteration adds a single application of $P_\Free^{-1}$
\cite{golub2013,benzi2005}. Proposition~\ref{prop:cg} then holds with $\kappa$
replaced by $\kappa_P = \kappa(P_\Free^{-1}H_{\Free\Free})$, so a preconditioner helps
exactly when $\kappa_P < \kappa(H_{\Free\Free})$, and is ideal when $P_\Free =
H_{\Free\Free}$.
\index{preconditioning}

\subsection{Preconditioning on a \emph{changing} free set}

Here the active-set setting imposes a constraint that a general treatment of
preconditioning does not, and it sorts the standard choices sharply. The free set
changes between outer steps, so a preconditioner is cheap only if it \emph{restricts}
to the new free set at negligible cost.

\paragraph{Jacobi restricts exactly.} $P = \diag(H)$ satisfies
$(\diag H)_\Free = \diag(H_\Free)$, so $P_\Free$ is the free-set slice of a single
vector computed once. For $H = M\T M$ that vector is the squared column norms of $M$,
available without forming $H$, and its application is $O(|\Free|)$. It carries no
per-step setup and is valid on every free set the loop visits---the natural default
when $H$ has a widely varying diagonal.

\paragraph{Incomplete Cholesky does not restrict.} In general
$(P_\Free)^{-1} \ne (P^{-1})_\Free$, and the factor of the full $P$ is not a factor of
$P_\Free$, so honouring $H_{\Free\Free}$ would require refactorising on each free set,
reintroducing exactly the assembly the method avoids. Such preconditioners pay off here
only once the free set stabilises---as it does near convergence, and across the
support-stable warm starts of Proposition~\ref{prop:warm}---where one factor serves many
outer steps.

\paragraph{Low-rank-plus-diagonal sits between.} $P = \diag(d) + U\Delta U\T$ built
from a few dominant eigenpairs \cite{saad2003iterative} restricts consistently in its
\emph{application}, $P_\Free = \diag(d)_\Free + U_\Free\Delta U_\Free\T$, and is
inverted matrix-free in $O(|\Free|r^2 + r^3)$ by Woodbury---precisely the direct
free-block solve of Chapter~\ref{ch:operator}. It captures the leading spectral
structure a diagonal misses. And when $H$ is \emph{itself} diagonal-plus-low-rank the
preconditioner \emph{is} $H$, giving $\kappa_P = 1$ and a direct inner solve that
bypasses CG altogether. Randomised Nyström sketching is the practical way to build one,
in a resketch-per-block variant and a sketch-once-and-mask variant that trade setup
cost against reuse.

\begin{remark}[The stopping criterion must not drift]
PCG naturally monitors the \emph{preconditioned} residual norm
$(r\T P_\Free^{-1}r)^{1/2}$, whereas Lemma~\ref{lem:inexact}---the transfer of finite
termination to inexact solves---is stated for the ordinary residual. The two differ by
a factor between $\lambda_{\min}(P_\Free)^{1/2}$ and $\lambda_{\max}(P_\Free)^{1/2}$.
To keep the guarantee intact, evaluate the inner stop on the unpreconditioned residual
(or tighten the preconditioned tolerance by $\kappa(P_\Free)^{1/2}$). The lemma then
applies word for word, because preconditioning changes only how fast the residual falls,
leaving the operator, the residual itself, and the decision margin untouched.

This is a small point with a general moral: when two guarantees are composed, check
that they are stated in the same units.
\end{remark}

\begin{remark}[Preconditioning widens the Krylov advantage]
A scaled gradient step minimises $f$ in the $P$-metric, but the orthant and simplex
projections that make projected gradient cheap are closed-form only in the Euclidean
norm, and become non-negative quadratic programs under a general $P$. So
preconditioning is a lever the Krylov inner solver has that the first-order reference
effectively lacks. It \emph{widens} the $\sqrt\kappa$ gap of
Chapter~\ref{ch:krylov} rather than closing it.
\end{remark}

\section{Linear shrinkage, and conditioning as a by-product}
\label{sec:shrinkage}

Now the statistical instance, which is where the chapter's thesis lives.

The sample covariance $\hat\Sig = X\T X/T$ is a poor estimator of $\Sig$ when $n$ is
not small relative to $T$: its smallest eigenvalues are severely downward-biased and
its condition number far exceeds that of $\Sig$. Ledoit and Wolf \cite{ledoit2004}
propose \emph{linear shrinkage}---replace $\hat\Sig$ by a convex combination with an
SPD target,
\[
  \hat\Sig_{\mathrm{LW}} = (1-\alpha)\,\hat\Sig + \alpha\, R\T R,
\]
which is exactly~\eqref{eq:regsplit}. The intensity $\alpha$ is chosen analytically by
minimising a consistent estimator of the expected \emph{Frobenius} loss; the classical
target is the scaled identity $\bar\lambda I$ with
$\bar\lambda = \tr(\hat\Sig)/n$.
\index{Ledoit--Wolf shrinkage}

Here is the tension, and its resolution.

\paragraph{The statistically optimal intensity barely helps conditioning.} On S\&P~500
data with $n/T \approx 0.41$, the Ledoit--Wolf oracle gives
$\alpha_{\mathrm{LW}} \approx 0.017$, and the Oracle Approximating Shrinkage
estimator \cite{chen2010} gives $\approx 0.010$. Both are small, and the mechanism is
structural rather than accidental. The oracle is $\alpha^\star = \beta^2/\delta^2$ with
$\delta^2 = \norm{\hat\Sig - \bar\lambda I}_F^2$, and $\delta^2$ is \emph{dominated by
the large, well-estimated eigenvalues}: the market factor alone contributes
$(\lambda_1 - \bar\lambda)^2 \approx (0.067)^2$, while each near-zero noise eigenvalue
contributes only $\bar\lambda^2 \approx 2\times10^{-7}$. The near-zero eigenvalues---the
sole cause of $\kappa \approx 109{,}000$---receive negligible correction, because the
Frobenius loss barely notices them. Corollary to Proposition~\ref{prop:weyl}: at
$\alpha^\star \approx 0.017$ the bound gives $\kappa \lesssim 8{,}500$, barely below
the unshrunk value.

\paragraph{The intensity practitioners actually use does help, decisively.}
Practitioners routinely apply $\alpha \approx 0.3$--$0.5$ for out-of-sample stability
\cite{demiguel2009,ledoit2003jpm}. At $\alpha = 0.5$ the bound gives
$\kappa \le \lambda_{\max}(\hat\Sig)/\bar\lambda + 1 \approx 147$, which matches the
empirical value---the bound is near-tight because
$\lambda_{\min}(\hat\Sig) \ll \alpha\bar\lambda$, so
$\lambda_{\min}(\hat\Sig_{\mathrm{LW}}) \approx \alpha\bar\lambda$, precisely the
denominator Weyl's inequality assumes. The measured effect on the solver is
proportionate: CG iterations fall from $684$ to $82$, and the outer active-set loop
contracts from $18$ steps to $6$, because a lower condition number moves the
unconstrained portfolio toward equal weights so fewer assets receive negative weights
in the primal step.

\paragraph{Is there a statistics-versus-speed trade-off?} The Frobenius surrogate and
the condition number are minimised at different $\alpha$, which looks like one. A
rolling-window out-of-sample study says there is none: realised out-of-sample variance
is minimised at $\alpha \approx 0.3$, inside the heavy-shrinkage range that gives the
conditioning benefit, while turnover falls monotonically in $\alpha$. \emph{The
intensity that makes the solver fast is also the intensity a risk-conscious
practitioner would choose}, and the apparent tension is an artefact of the Frobenius
surrogate rather than a property of the problem.

\paragraph{Factor targets.} The scaled identity is not the only choice. A structured
target $T_0 = BB\T + D$, with $B$ holding the loadings of $k \ll n$ risk factors and
$D$ a diagonal idiosyncratic variance, is a drop-in replacement in the stacked-matrix
construction: take $R = (B\T; D^{1/2})$. Because the factor structure clusters
eigenvalues more tightly than a scaled identity, Proposition~\ref{prop:weyl} yields a
smaller bound at the same $\alpha$. The product cost is unchanged, the factor and
diagonal terms costing $O(nk)$ and $O(n)$ against the data term's $O(nT)$.

\section{Spectral cleaning: Marchenko--Pastur}

Shrinkage moves every eigenvalue toward the target. Random-matrix theory offers
something sharper: identify \emph{which} eigenvalues are noise and replace only those.
\index{Marchenko--Pastur law}

Work with the sample \emph{correlation} $C = D^{-1/2}\hat\Sig D^{-1/2}$, where
$D = \diag(\hat\Sig)$. This matters: individual variances are well estimated even in
high dimensions, and the noise that destabilises the problem lives in the correlation
structure. The covariance mixes the well-estimated variances into the spectrum, so its
bulk is not a clean band and its edge has no scale-free form.

The Marchenko--Pastur law \cite{marchenko1967} describes the limiting spectral
distribution of a sample correlation of standardised series when $n, T \to \infty$
with $n/T \to c \in (0,1)$: under the null of no correlation structure the bulk is
supported on $[\lambda_-,\lambda_+]$ with $\lambda_\pm = (1 \pm \sqrt{c})^2$. In
finite samples the empirical upper edge is $\hat\lambda_+ = (1 + \sqrt{n/T})^2$.
Correlation eigenvalues above it carry genuine factor structure; those below are
consistent with noise. Write $k$ for the number above. On S\&P~500 data with $n = 494$,
$T = 1213$: $\hat\lambda_+ \approx 2.28$ and $k = 17$.

\subsection{The constant-residual-eigenvalue estimator}

Let $C = U\Lambda U\T$, partition by the edge into signal $U_k, \Lambda_k$ and noise
$U_0$, and replace the bulk eigenvalues by their average
\[
  \bar\mu = \frac{1}{n-k}\sum_{j>k}\lambda_j = \frac{n - \sum_{j\le k}\lambda_j}{n-k},
\]
which preserves the trace $\tr(C) = n$. The resolution of the identity
$U_kU_k\T + U_0U_0\T = I$ then gives the \emph{constant residual eigenvalue} (CRE), or
eigenvalue-clipping, cleaned correlation \cite{laloux1999,bun2017,lopezdeprado2018}
\begin{equation}
  C_0 = U_k\Lambda_k U_k\T + \bar\mu\,U_0U_0\T = \bar\mu\,I + U_k\Delta_k U_k\T,
  \qquad \Delta_k = \Lambda_k - \bar\mu I,
  \label{eq:cre}
\end{equation}
which is \emph{scalar-identity-plus-low-rank}, with $\delta_j = \lambda_j - \bar\mu > 0$.

\begin{proposition}[Spectral properties of $C_0$]
\label{prop:rmt_spectrum}
$C_0 \succ 0$ with eigenvalues $\{\lambda_1,\dots,\lambda_k,\bar\mu,\dots,\bar\mu\}$
and $\tr(C_0) = n$; $\kappa(C_0) = \lambda_{\max}(C)/\bar\mu$; $C_0$ differs from $C$
only in the noise subspace; and $C_0 v = \bar\mu v + U_k(\Delta_k(U_k\T v))$ costs
$O(nk)$ with $O(nk)$ storage.
\end{proposition}

\begin{proof}
All four from the eigendecomposition
$C_0 = [U_k\ U_0]\diag(\Lambda_k,\bar\mu I)[U_k\ U_0]\T$, the positivity of
$\delta_j$, and the definition of $\bar\mu$.
\end{proof}

Cleaning lifts the noise floor to $\bar\mu$ while preserving the signal, so
$\kappa(C_0) \approx 313$ against $\kappa(C) \approx 91{,}000$: an improvement of
roughly $290\times$, and decisively better than what heavy shrinkage achieves.

\subsection{Two design decisions worth understanding}

\begin{remark}[Why solve in standardised coordinates]
\label{rem:standardised}
Solve the problem in $y = D^{1/2}w$, posing it against $C_0$ rather than rebuilding the
covariance $\Sig_0 = D^{1/2}C_0D^{1/2}$. The reason is arithmetic, not aesthetic: the
rebuilt covariance is \emph{far worse conditioned} ($\kappa(\Sig_0) \approx 1{,}400$ on
the same data), inflated by the spread of asset variances---a factor of about $49$
between largest and smallest. The change of variables is exact, and it also puts the
tolerance on an $O(1)$ scale with noise floor $\bar\mu \approx 0.5$, removing the
dependence on the raw data scale that a covariance formulation carries.
\end{remark}

\begin{remark}[Why clipping, and not a better estimator]
\label{rem:rie}
CRE clipping is the simplest rotationally invariant estimator and is \emph{not}
loss-optimal: nonlinear shrinkage \cite{ledoit2012,ledoit2020} and the oracle RIE
\cite{ledoitpeche2011,bun2017}, which reshape \emph{every} eigenvalue, dominate it in
Frobenius loss. We nevertheless use it for a \emph{structural} reason: clipping is the
only member of the family whose cleaned correlation is
scalar-identity-plus-low-rank---the spiked form \cite{johnstone2001} that admits the
closed-form Woodbury solve of Chapter~\ref{ch:operator}. A full-rank RIE modification
would require a general factorisation and forfeit it.

Whether the statistical gap matters for the application is an empirical question, and
on long-only minimum-variance data it is small. But note the shape of the reasoning:
a slightly worse estimator that is $O(nk)$ to invert may dominate a slightly better one
that is $O(n^3)$, once one is computing hundreds of paths rather than one point.
\emph{Statistical optimality and computational structure are both objectives, and they
are traded, not ranked.}
\end{remark}

Two further notes. A blend $(1-\alpha)C + \alpha C_0$ would re-introduce a fraction of
the bulk eigenvalues, and---the algebraic point---only at $\alpha = 1$ does the
restricted system retain the scalar-identity-plus-low-rank form; for
$\alpha \in (0,1)$ the residual structure is full rank and no $k\times k$ correction
exists. And the solver applies for \emph{any} choice of $k$: a practitioner with an
external risk model, or a PCA with $k$ fixed by cross-validation, can supply the
eigenpairs directly. What the solve requires is not the source of $k$ but the
replacement of the remaining $n-k$ eigenvalues by a single scalar.

\subsection{Is $k$ a delicate hyperparameter?}

No, and it is worth knowing why one can say so. Two biases perturb the threshold:
finite-$T$ Tracy--Widom fluctuations of the largest noise eigenvalue, at scale
$O(T^{-2/3})$ \cite{johnstone2001,elkaroui2008}, and diffuse sector structure just below
the edge. Both move $k$ by at most an eigenvalue or two. An audit at $k-1$, $k$, $k+1$
costs one extra solve per audit point and reveals the sensitivity directly, which is
small over the bulk of a frontier.

That is the right way to treat a threshold in general: rather than argue that it is
correct, measure what changes when it moves by the amount it could plausibly be wrong
by.

\begin{notes}
Ledoit and Wolf's linear shrinkage is \cite{ledoit2004}, with OAS \cite{chen2010} and
nonlinear shrinkage \cite{ledoit2020}. Section~\ref{sec:shrinkage} follows
\cite{schmelzer2026matrixfree}, whose contribution is the identification of the
shrinkage intensity as a preconditioning parameter through the $\alpha$-dependent
bound, and the out-of-sample study reconciling the two objectives. The RMT material is
\cite{schmelzer2026rmt}; Marchenko--Pastur is \cite{marchenko1967}, eigenvalue clipping
\cite{laloux1999}, and the modern treatment \cite{bun2017}. Preconditioning on a
changing free set, including the Nyström constructions, is from
\cite{schmelzer2026nncg}; Benzi \cite{benzi2005} and Saad \cite{saad2003iterative} are
the general references.
\end{notes}

\begin{exercises}

\begin{exercise}
Prove Proposition~\ref{prop:weyl} from Weyl's inequality, and verify the
scaled-identity corollary.
\end{exercise}

\begin{exercise}
Verify that $\tilde M\T\tilde M = H_\alpha$ for the stacked matrix of Section~1, and
explain why this keeps the method matrix-free where forming $H_\alpha$ would not.
\end{exercise}

\begin{exercise}
Show that Jacobi preconditioning restricts exactly, $(\diag H)_\Free = \diag(H_\Free)$,
and that incomplete Cholesky does not, by exhibiting a $3\times3$ example.
\end{exercise}

\begin{exercise}
Given $\lambda_{\max}(\hat\Sig) = 0.07$, $\bar\lambda = 4.6\times10^{-4}$, and
$\lambda_{\min}(\hat\Sig) = 6\times10^{-7}$, evaluate the bound of
Proposition~\ref{prop:weyl} at $\alpha \in \{0.017, 0.1, 0.3, 0.5\}$ and confirm the
numbers quoted in Section~\ref{sec:shrinkage}.
\end{exercise}

\begin{exercise}
Explain, using $\delta^2 = \norm{\hat\Sig - \bar\lambda I}_F^2$, why the Ledoit--Wolf
oracle intensity is insensitive to the near-zero eigenvalues. Construct a spectrum for
which it would not be.
\end{exercise}

\begin{exercise}
Verify Proposition~\ref{prop:rmt_spectrum}, and compute $\bar\mu$ and $\kappa(C_0)$ for
a spectrum with $n = 100$, $k = 3$, $\lambda_{1,2,3} = (20, 8, 4)$ and the remaining
trace spread over the bulk.
\end{exercise}

\begin{exercise}
Show that a blend $(1-\alpha)C + \alpha C_0$ is scalar-identity-plus-low-rank only for
$\alpha \in \{0,1\}$, confirming Section~5.2.
\end{exercise}

\begin{exercise}[Implementation]
Download or simulate a $T\times n$ return matrix with $n/T \approx 0.4$. Compute
$\kappa$ of the sample covariance, of the correlation, of the CRE-cleaned correlation,
and of the rebuilt cleaned covariance. Confirm the ordering claimed in
Remark~\ref{rem:standardised}.
\end{exercise}

\begin{exercise}[Implementation]
Solve the long-only minimum-variance problem with matrix-free CG under
$\alpha \in \{0, 0.017, 0.1, 0.3, 0.5\}$ and record CG iterations and outer active-set
steps for each. Reproduce the qualitative shape of the $684 \to 82$ collapse, and
explain why the outer count falls too.
\end{exercise}

\end{exercises}
